% !TEX program = xelatex
% Compile with XeLaTeX. All content and layout are editable in this file.
% Exactly two slides. The weekly format is the Profilprojekt (IMO) adaptation.
% Source for the Scrum principles on slide 1:
% https://scrumguides.org/docs/scrumguide/v2020/2020-Scrum-Guide-German.pdf
% Weekly timing and the agenda on slide 2 are facilitation recommendations.
\documentclass[12pt,aspectratio=169]{beamer}

\usepackage{fontspec}
\usepackage[ngerman]{babel}
\usepackage[absolute,overlay]{textpos}
\setsansfont{texgyreheros-regular.otf}[
  BoldFont=texgyreheros-bold.otf,
  ItalicFont=texgyreheros-italic.otf,
  BoldItalicFont=texgyreheros-bolditalic.otf]
\setmainfont{texgyreheros-regular.otf}[
  BoldFont=texgyreheros-bold.otf,
  ItalicFont=texgyreheros-italic.otf,
  BoldItalicFont=texgyreheros-bolditalic.otf]
\usefonttheme{professionalfonts}

\geometry{paperwidth=320mm,paperheight=180mm}
\setlength{\TPHorizModule}{1mm}
\setlength{\TPVertModule}{1mm}
\setlength{\parindent}{0pt}

\definecolor{Paper}{HTML}{F4F2EB}
\definecolor{Ink}{HTML}{0B2F38}
\definecolor{Muted}{HTML}{50646A}
\definecolor{Accent}{HTML}{A66F08}
\definecolor{Line}{HTML}{C9D3D2}
\setbeamercolor{background canvas}{bg=Paper}
\setbeamercolor{normal text}{fg=Ink,bg=Paper}
\setbeamertemplate{navigation symbols}{}
\setbeamertemplate{headline}{}
\setbeamertemplate{footline}{}
\setbeamertemplate{frametitle}{}
\title{Weekly Scrum im Profilprojekt (IMO)}
\author{Prof. Dr. Raphael Volz (HSPF)}
\date{2. November 2026}
\hypersetup{
  pdftitle={Weekly Scrum im Profilprojekt (IMO)},
  pdfauthor={Prof. Dr. Raphael Volz (HSPF)},
  pdfsubject={Leitfaden und Scrum-Prinzipien, 2. November 2026},
  pdfkeywords={Weekly Scrum, Standup, Profilprojekt, IMO, HSPF},
  colorlinks=true,urlcolor=Muted,
  pdfdisplaydoctitle=true
}

% Absolute placement in millimetres on a 320 x 180 mm slide.
% Arguments: x, y, width, font size, line spacing, color, content.
\newcommand{\Place}[7]{%
  \begin{textblock}{#3}(#1,#2)
    \raggedright\color{#6}\fontsize{#4}{#5}\selectfont #7\par
  \end{textblock}%
}
\newcommand{\SlideHeader}[2]{%
  \Place{18}{11}{284}{12}{14}{Muted}{Profilprojekt (IMO)}%
  \Place{18}{22}{284}{32}{36}{Ink}{\bfseries #1}%
  \Place{18}{40}{284}{19}{24}{Muted}{#2}%
}
\newcommand{\SlideFooter}[1]{%
  \Place{18}{165}{284}{9}{11}{Line}{\rule{284mm}{0.35pt}}%
  \Place{18}{169}{173}{10}{12}{Muted}{Prof. Dr. Raphael Volz (HSPF)}%
  \Place{203}{169}{80}{10}{12}{Muted}{2. November 2026}%
  \Place{293}{169}{9}{10}{12}{Muted}{\hfill #1/2}%
}

\begin{document}

% SLIDE 1 -- Scrum principles, applied to this project's weekly format.
% Deliberately does not present "Weekly Scrum" as a prescribed Scrum event.
\begin{frame}[plain]
  \SlideHeader{Weekly Scrum im Profilprojekt (IMO)}%
    {\textbf{Scrum} ist ein agiles Rahmenwerk für die schrittweise Lösung komplexer Aufgaben.\\
     Teams arbeiten in kurzen Zyklen (Sprints), prüfen Ergebnisse und passen ihr Vorgehen an.}

  \Place{18}{61}{284}{13}{17}{Accent}{\bfseries DREI SCRUM-PRINZIPIEN IM WÖCHENTLICHEN PROJEKTRHYTHMUS}

  \Place{18}{74}{86}{24}{29}{Ink}{\bfseries Transparenz}
  \Place{18}{90}{86}{19}{25}{Ink}{Ergebnisse, offene Aufgaben und Hindernisse im gemeinsamen Board sichtbar machen.}

  \Place{118}{74}{86}{24}{29}{Ink}{\bfseries Überprüfung}
  \Place{118}{90}{86}{19}{25}{Ink}{Den Fortschritt zum Projektziel prüfen und Abweichungen früh erkennen.}

  \Place{218}{74}{84}{24}{29}{Ink}{\bfseries Anpassung}
  \Place{218}{90}{84}{19}{25}{Ink}{Arbeit priorisieren, Unterstützung vereinbaren und die nächsten Schritte festlegen.}

  \Place{18}{130}{284}{19}{25}{Ink}{\bfseries Selbstorganisation: Alle Teammitglieder stimmen ihre Arbeit gemeinsam ab.}
  \Place{18}{143}{284}{18}{24}{Ink}{Ergebnis: ein gemeinsamer Wochenplan mit Zuständigkeiten und Terminen.}
  \Place{18}{156}{284}{10}{13}{Muted}{Projektformat: wöchentlich angepasst. Grundlage: \href{https://scrumguides.org/docs/scrumguide/v2020/2020-Scrum-Guide-German.pdf}{Schwaber \& Sutherland, Scrum Guide (2020), Definition, Theorie und Team, S. 3--5}.}
  \SlideFooter{1}
\end{frame}

% SLIDE 2 -- Reusable, projected facilitation guide.
\begin{frame}[plain]
  \SlideHeader{Weekly Scrum: Leitfaden für das Standup}%
    {Gemeinsam Ergebnisse prüfen, die nächste Woche planen und Unterstützung klären}

  \Place{18}{59}{12}{24}{28}{Accent}{\bfseries 01}
  \Place{36}{59}{181}{22}{27}{Ink}{\bfseries Was ist seit letzter Woche fertig?}
  \Place{36}{72}{177}{18}{23}{Muted}{Zeigen Sie ein Ergebnis und benennen Sie offene Aufgaben.}

  \Place{18}{88}{12}{24}{28}{Accent}{\bfseries 02}
  \Place{36}{88}{181}{22}{27}{Ink}{\bfseries Was erreichen Sie bis nächste Woche?}
  \Place{36}{101}{177}{18}{23}{Muted}{Nennen Sie Ihren nächsten Schritt und das erwartete Ergebnis.}

  \Place{18}{117}{12}{24}{28}{Accent}{\bfseries 03}
  \Place{36}{117}{181}{22}{27}{Ink}{\bfseries Wo brauchen Sie Unterstützung?}
  \Place{36}{130}{177}{18}{23}{Muted}{Machen Sie Hindernisse und Abhängigkeiten sichtbar.}

  \Place{234}{59}{68}{22}{27}{Ink}{\bfseries 15 Minuten}
  \Place{234}{71}{68}{12}{16}{Muted}{Ablaufvorschlag pro Team}
  \Place{234}{87}{23}{17}{22}{Accent}{\bfseries 2 Min.}
  \Place{259}{87}{43}{17}{22}{Ink}{Ziel \& Board}
  \Place{234}{100}{23}{17}{22}{Accent}{\bfseries 8 Min.}
  \Place{259}{100}{43}{17}{22}{Ink}{Teamrunde}
  \Place{234}{113}{23}{17}{22}{Accent}{\bfseries 3 Min.}
  \Place{259}{113}{43}{17}{22}{Ink}{Hilfe klären}
  \Place{234}{126}{23}{17}{22}{Accent}{\bfseries 2 Min.}
  \Place{259}{126}{43}{17}{22}{Ink}{Abschluss}

  \Place{18}{147}{284}{17}{22}{Ink}{\bfseries Im Board festhalten: nächster Schritt, verantwortliche Person und Termin.}
  \Place{18}{157}{284}{14}{18}{Muted}{Details klären Sie im Anschluss. Hindernisse sprechen Sie auch zwischen den Treffen an.}
  \SlideFooter{2}
\end{frame}

\end{document}
